\section{Familias de Optimizadores}
\label{sec:annex-optimizers}

\subsection{Comparación de Memoria y Complejidad}

La Tabla~\ref{tab:optimizer-complexity} resume la sobrecarga de memoria (número de búferes de estado por parámetro), la complejidad computacional por paso y los hiperparámetros clave para todos los optimizadores soportados. La sobrecarga de memoria se expresa en términos del número de tensores de estado mantenidos por parámetro del modelo, donde cada tensor tiene la misma forma que el parámetro.

\begin{longtable}{llcl}
\toprule
\textbf{Optimizador} & \textbf{Búferes de Estado} & \textbf{Costo por Paso} & \textbf{Hiperparámetros Clave} \\
\midrule
\endfirsthead
\toprule
\textbf{Optimizador} & \textbf{Búferes de Estado} & \textbf{Costo por Paso} & \textbf{Hiperparámetros Clave} \\
\midrule
\endhead
\midrule
\multicolumn{4}{r}{\textit{continúa en la página siguiente}} \\
\endfoot
\bottomrule
\caption{Comparación de memoria y complejidad de todos los optimizadores soportados. $n$ = número de parámetros, $d$ = dimensionalidad del tensor, $r$ = dimensión de rango bajo, $k$ = número de iteraciones de Newton--Schulz. Búferes de estado indica el número de tensores de estado del optimizador mantenidos por grupo de parámetros. El costo por paso se enfoca en el costo adicional dominante más allá del propio cálculo del gradiente.}
\label{tab:optimizer-complexity}
\endlastfoot
SGD+Momentum & 1 ($m$) & $\mathcal{O}(n)$ & lr, momentum, wd \\
AdamW & 2 ($m$, $v$) & $\mathcal{O}(n)$ & lr, $\beta_1$, $\beta_2$, eps, wd \\
RAdam & 2 ($m$, $v$) & $\mathcal{O}(n)$ & lr, $\beta_1$, $\beta_2$, eps, wd \\
Adan & 3 ($m$, $v$, $s$) & $\mathcal{O}(n)$ & lr, $\beta_1$, $\beta_2$, $\beta_3$, eps, wd \\
ADOPT & 2 ($m$, $v$) & $\mathcal{O}(n)$ & lr, $\beta_1$, $\beta_2$, eps, wd \\
AdEMAMix & 3 ($m_1$, $m_2$, $v$) & $\mathcal{O}(n)$ & lr, $\beta_1$, $\beta_2$, $\beta_3$, eps, wd \\
MARS & 3 ($m$, $v$, $z$) & $\mathcal{O}(n)$ & lr, $\beta_1$, $\beta_2$, eps, wd, $\gamma$ \\
Cautious AdamW & 2 ($m$, $v$) & $\mathcal{O}(n)$ & lr, $\beta_1$, $\beta_2$, eps, wd \\
LAMB & 2 ($m$, $v$) & $\mathcal{O}(n)$ & lr, $\beta_1$, $\beta_2$, eps, wd \\
Schedule-Free & 3 ($z$, $x$, $n$) & $\mathcal{O}(n)$ & lr, $\beta_1$, $\beta_2$, wd \\
Adafactor & 1--2 (row/col) & $\mathcal{O}(n)$ & lr, $\beta_1$, $\beta_2$, eps, wd, factored \\
GaLore & 2 ($m$, $v$) + SVD/proj & $\mathcal{O}(nr)$ & lr, rank $r$, $\beta_1$, $\beta_2$, eps, wd \\
Prodigy & 3 ($m$, $v$, $d$) & $\mathcal{O}(n)$ & $\beta_1$, $\beta_2$, eps, wd, $d_0$ \\
Lion & 1 ($m$) & $\mathcal{O}(n)$ & lr, $\beta_1$, $\beta_2$, wd \\
Shampoo & $2d$ ($L_i$, $R_i$) & $\mathcal{O}(n^{1+2/d})$ & lr, eps, wd, matrix eps \\
SOAP & $2d$ + 2 ($m$, $v$) & $\mathcal{O}(n^{1+1/d})$ & lr, $\beta_1$, $\beta_2$, eps, wd, shampoo eps \\
Sophia & 3 ($m$, $v$, $h$) & $\mathcal{O}(n)$ & lr, $\beta_1$, $\beta_2$, eps, wd, $k$ \\
Muon & 2 ($m$, $v$) & $\mathcal{O}(n \cdot k)$ & lr, momentum, wd, NS steps $k$ \\
Turbo-Muon & 2 ($m$, $v$) & $\mathcal{O}(n \cdot k)$ & lr, momentum, wd, NS steps $k$ \\
APOLLO & 2 bajo rango ($m_R$, $v_R$) + proy & $\mathcal{O}(nr)$ & lr, rank $r$, update gap, scale, betas, eps, wd \\
APOLLO-Mini & 2 rango-1 ($m_R$, $v_R$) + proy & $\mathcal{O}(n)$ & lr, update gap, scale, betas, eps, wd \\
Q-APOLLO & 2 bajo rango cuantizados + proy & $\mathcal{O}(nr)$ & lr, rank $r$, update gap, scale, quant bits, betas, eps, wd \\
Anon & 3 ($m$, $v$, fixed\_lr) & $\mathcal{O}(n)$ & lr, $\beta_1$, $\beta_2$, eps, wd, $\gamma$ \\
\end{longtable}

\subsection{La Evolución de la Optimización en Redes Neuronales}
El estudio de optimizadores enmarca la optimización de transformadores como una respuesta a tres presiones estructurales: paisajes de pérdida no convexos, heterogeneidad severa de curvatura entre bloques de parámetros y el costo de memoria de almacenar el estado del optimizador para modelos muy grandes. El informe sostiene que el campo se ha diversificado en varias trayectorias: líneas base adaptativas de primer orden, métodos de reducción de varianza, métodos eficientes en memoria, precondicionadores estructurados de segundo orden, métodos sin programación de tasa de aprendizaje y actualizaciones orientadas a ortogonalidad.

\subsection{Línea Base Estándar y Optimizadores Adaptativos}
\paragraph{SGD con Momentum.}
La actualización clásica acumula un búfer de momentum y luego aplica una tasa de aprendizaje fija. En cada iteración $t$, el algoritmo mantiene un promedio móvil exponencial $m_t$ de los gradientes pasados:
\begin{align}
m_t &= \beta m_{t-1} + g_t \\
\theta_{t+1} &= \theta_t - \eta m_t
\end{align}
donde $g_t = \nabla f(\theta_t)$, $\beta \in [0.8, 0.99]$ controla la decadencia del momentum y $\eta$ es la tasa de aprendizaje fija. El término de momentum actúa como velocidad: acelera el movimiento en direcciones consistentes mientras amortigua oscilaciones en direcciones variables. Sus fortalezas son la baja sobrecarga de memoria (un único búfer de momentum por parámetro) y una fuerte generalización cuando se ajusta cuidadosamente. Su principal debilidad en cargas de trabajo de transformadores es la escasa robustez frente a la heterogeneidad de curvatura (diferentes espectros del Hessiano entre grupos de parámetros) y una fuerte dependencia de las programaciones de tasa de aprendizaje.

\begin{algorithm}[H]
\caption{SGD con Momentum}
\label{alg:sgd-momentum}
\begin{algorithmic}[1]
\Require Parámetros iniciales $\theta_0$, tasa de aprendizaje $\eta$, coeficiente de momentum $\beta$, decadencia de peso $\lambda$
\Ensure Parámetros actualizados $\theta_T$
\State Inicializar: búfer de momentum $m \gets 0$
\For{$t = 0, 1, 2, \ldots, T-1$}
    \State Calcular gradiente: $g_t \gets \nabla f(\theta_t)$
    \If{decadencia\_de\_peso $> 0$}
        \State $g_t \gets g_t + \lambda \theta_t$ \Comment{Regularización L2}
    \EndIf
    \State Actualizar momentum: $m \gets \beta \cdot m + g_t$
    \State Actualizar parámetros: $\theta_{t+1} \gets \theta_t - \eta \cdot m$
\EndFor
\Return $\theta_T$
\end{algorithmic}
\end{algorithm}

\paragraph{Adam y AdamW.}
Adam rastrea promedios móviles exponenciales tanto del primer momento (media) como del segundo momento (varianza no centrada) para lograr tasas de aprendizaje adaptativas elemento a elemento. La regla de actualización es:
\begin{align}
m_t &= \beta_1 m_{t-1} + (1-\beta_1) g_t \\
v_t &= \beta_2 v_{t-1} + (1-\beta_2) g_t^2 \\
\hat{m}_t &= \frac{m_t}{1-\beta_1^t}, \quad \hat{v}_t = \frac{v_t}{1-\beta_2^t} \quad \text{(corrección de sesgo)} \\
\theta_{t+1} &= \theta_t - \eta \left(\frac{\hat{m}_t}{\sqrt{\hat{v}_t}+\epsilon}\right)
\end{align}
AdamW introduce una modificación crucial: la decadencia de peso se aplica directamente a los parámetros (desacoplada de los gradientes): $\theta_t \gets \theta_t(1-\eta\lambda)$ en lugar de añadir $\lambda\theta_t$ al gradiente. Este desacoplamiento evita que el escalado adaptativo interfiera con la fuerza de la regularización. El informe trata a AdamW como la línea base práctica para el ajuste fino de transformadores porque converge rápidamente y es relativamente tolerante a la variación de hiperparámetros. La desventaja es el costo de memoria: ambos tensores de momento deben almacenarse para cada parámetro, duplicando la memoria del estado del optimizador en relación con SGD.

\begin{algorithm}[H]
\caption{AdamW (Adam con Decadencia de Peso Desacoplada)}
\label{alg:adamw}
\begin{algorithmic}[1]
\Require Parámetros iniciales $\theta_0$, tasa de aprendizaje $\eta$, tasas de decadencia exponencial $\beta_1, \beta_2 \in [0,1)$
\Require Constante de momentum $\epsilon$, decadencia de peso $\lambda$
\Ensure Parámetros actualizados $\theta_T$
\State Inicializar: primer momento $m \gets 0$, segundo momento $v \gets 0$, contador de pasos $t \gets 0$
\For{$t = 1, 2, \ldots, T$}
    \State Calcular gradiente: $g_t \gets \nabla f(\theta_{t-1})$
    \State Decadencia de peso (desacoplada): $\theta_{t-1} \gets \theta_{t-1}(1 - \eta\lambda)$
    \State Actualizar momentos: $m \gets \beta_1 m + (1-\beta_1) g_t$
    \State $v \gets \beta_2 v + (1-\beta_2) g_t^2$
    \State Corrección de sesgo: $\hat{m} \gets m / (1 - \beta_1^t)$, $\hat{v} \gets v / (1 - \beta_2^t)$
    \State Actualizar parámetros: $\theta_t \gets \theta_{t-1} - \eta (\hat{m} / (\sqrt{\hat{v}} + \epsilon))$
\EndFor
\Return $\theta_T$
\end{algorithmic}
\end{algorithm}

\paragraph{RAdam (Adam Rectificado).}
RAdam aborda la inestabilidad de Adam en el entrenamiento temprano rectificando dinámicamente la tasa de aprendizaje adaptativa. La observación clave es que el segundo momento $v_t$ tiene una varianza muy alta en los primeros pasos, causando un escalado adaptativo poco fiable. RAdam calcula la longitud efectiva de la ventana de promedio móvil simple (SMA):
\[\rho_t = \rho_{\infty} - \frac{2t\beta_2^t}{1-\beta_2^t}, \quad \text{donde} \quad \rho_{\infty} = \frac{2}{1-\beta_2}-1\]
Cuando $\rho_t > 4$ (suficientes muestras para la estimación de varianza), RAdam aplica el escalado adaptativo con un término de rectificación:
\[r_t = \sqrt{\frac{(\rho_t-4)(\rho_t-2)\rho_{\infty}}{(\rho_{\infty}-4)(\rho_{\infty}-2)\rho_t}}, \quad \theta_{t+1} = \theta_t - \eta r_t \frac{\hat{m}_t}{\sqrt{\hat{v}_t}+\epsilon}\]
Cuando $\rho_t \leq 4$, RAdam retrocede a SGD con momentum: $\theta_{t+1} = \theta_t - \eta\hat{m}_t$. Esta transición gradual elimina la necesidad de programaciones manuales de calentamiento de la tasa de aprendizaje y mejora la robustez.

\begin{algorithm}[H]
\caption{RAdam (Adam Rectificado)}
\label{alg:radam}
\begin{algorithmic}[1]
\Require Parámetros iniciales $\theta_0$, tasa de aprendizaje $\eta$, $\beta_1, \beta_2$, decadencia de peso $\lambda$
\Ensure Parámetros actualizados $\theta_T$
\State Inicializar: $m \gets 0$, $v \gets 0$, $t \gets 0$
\State Calcular: $\rho_\infty \gets 2/(1-\beta_2) - 1$
\For{$t = 1, 2, \ldots, T$}
    \State $g_t \gets \nabla f(\theta_{t-1})$
    \If{decadencia\_de\_peso $> 0$}
        \State $\theta_{t-1} \gets \theta_{t-1}(1 - \eta\lambda)$
    \EndIf
    \State $m \gets \beta_1 m + (1-\beta_1) g_t$
    \State $v \gets \beta_2 v + (1-\beta_2) g_t^2$
    \State $\hat{m} \gets m/(1-\beta_1^t)$, $\hat{v} \gets v/(1-\beta_2^t)$
    \State $\rho_t \gets \rho_\infty - 2t \beta_2^t / (1 - \beta_2^t)$
    \If{$\rho_t > 4$}
        \State $r_t \gets \sqrt{((\rho_t-4)(\rho_t-2)\rho_\infty) / ((\rho_\infty-4)(\rho_\infty-2)\rho_t)}$
        \State $\theta_t \gets \theta_{t-1} - \eta r_t \hat{m} / (\sqrt{\hat{v}} + \epsilon)$
    \Else
        \State $\theta_t \gets \theta_{t-1} - \eta \hat{m}$ \Comment{SGD con momentum}
    \EndIf
\EndFor
\Return $\theta_T$
\end{algorithmic}
\end{algorithm}

\subsection{Momentum Avanzado y Reducción de Varianza (2024--2025)}
\paragraph{Adan (Momentum Nesterov Adaptativo).}
Adan reformula la aceleración de Nesterov sin el cálculo de gradiente adicional requerido por el SGD Nesterov clásico. El algoritmo mantiene tres búferes de momentum:
\begin{align}
m_t &= (1-\beta_1)m_{t-1} + \beta_1 g_t \quad\text{(primer momento)} \\
v_t &= (1-\beta_2)v_{t-1} + \beta_2(g_t - g_{t-1}) \quad\text{(velocidad/diferencia de gradientes)} \\
n_t &= (1-\beta_3)n_{t-1} + \beta_3[g_t + (1-\beta_1)(g_t-g_{t-1})]^2 \quad\text{(segundo momento Nesterov)}
\end{align}
El término de \textbf{Estimación de Momentum Nesterov} (NME) $\bar{g}_t = g_t + (1-\beta_1)(g_t-g_{t-1})$ estima el gradiente en una posición futura sin evaluarlo. La actualización combina el momentum con la velocidad:
\[\bar{m}_t = m_t + (1-\beta_1)v_t, \quad \theta_{t+1} = \theta_t - \eta\frac{\bar{m}_t}{\sqrt{n_t}+\epsilon}\]
Adan logra una convergencia rápida en diversas arquitecturas (CNN, GAN, Transformadores) mediante esta aceleración. El costo es mantener tres búferes similares al momentum, aumentando la sobrecarga de memoria en relación con Adam.

\begin{algorithm}[H]
\caption{Adan (Momentum Nesterov Adaptativo)}
\label{alg:adan}
\begin{algorithmic}[1]
\Require Parámetros iniciales $\theta_0$, tasa de aprendizaje $\eta$, $\beta_1, \beta_2, \beta_3 \in (0,1)$
\Ensure Parámetros actualizados $\theta_T$
\State Inicializar: $m \gets 0$, $v \gets 0$, $n \gets 0$, $g_{-1} \gets 0$ (gradiente anterior)
\For{$t = 1, 2, \ldots, T$}
    \State $g_t \gets \nabla f(\theta_{t-1})$
    \State $m \gets (1-\beta_1) m + \beta_1 g_t$
    \State $\Delta g_t \gets g_t - g_{t-1}$
    \State $v \gets (1-\beta_2) v + \beta_2 \Delta g_t$
    \State Estimación Nesterov: $\bar{g}_t \gets g_t + (1-\beta_1) \Delta g_t$
    \State $n \gets (1-\beta_3) n + \beta_3 \bar{g}_t^2$
    \State Momentum combinado: $\bar{m} \gets m + (1-\beta_1) v$
    \State $\theta_t \gets \theta_{t-1} - \eta (\bar{m} / (\sqrt{n} + \epsilon))$
    \State $g_{t-1} \gets g_t$
\EndFor
\Return $\theta_T$
\end{algorithmic}
\end{algorithm}

\paragraph{ADOPT (Adam con Ajuste Óptimo Podado).}
ADOPT corrige un problema teórico fundamental en Adam: el gradiente aparece tanto en la estimación del primer momento como en la del segundo momento, creando circularidad. ADOPT \textbf{desacopla} utilizando el segundo momento del paso anterior para el denominador:
\begin{align}
v_t &= \beta_2 v_{t-1} + (1-\beta_2)g_t^2 \\
m_t &= \beta_1 m_{t-1} + (1-\beta_1)\frac{g_t}{\sqrt{v_{t-1}}+\epsilon} \quad\text{(usa $v_{t-1}$, no $v_t$)} \\
\theta_{t+1} &= \theta_t - \eta m_t
\end{align}
Este simple reordenamiento logra la tasa de convergencia óptima $O(1/\sqrt{T})$ con \textbf{cualquier} elección de $\beta_2 \in (0,1)$, sin suposiciones de ruido acotado. La consecuencia práctica es que ADOPT es un reemplazo directo de Adam con garantías teóricas más sólidas y un rendimiento empírico comparable o superior en los dominios de visión, PLN, RL y modelado generativo.

\begin{algorithm}[H]
\caption{ADOPT (Adam con Ajuste Óptimo Podado)}
\label{alg:adopt}
\begin{algorithmic}[1]
\Require Parámetros iniciales $\theta_0$, tasa de aprendizaje $\eta$, $\beta_1, \beta_2$, tolerancia $\epsilon$
\Ensure Parámetros actualizados $\theta_T$
\State Inicializar: $m \gets 0$, $v \gets 0$
\For{$t = 1, 2, \ldots, T$}
    \State $g_t \gets \nabla f(\theta_{t-1})$
    \If{decadencia\_de\_peso $> 0$}
        \State $\theta_{t-1} \gets \theta_{t-1}(1 - \eta\lambda)$
    \EndIf
    \State Usar segundo momento \textbf{anterior}: $\text{denom} \gets \sqrt{\max(v, \epsilon)} + \epsilon$
    \State Actualizar primer momento usando varianza anterior: $m \gets \beta_1 m + (1-\beta_1)(g_t / \text{denom})$
    \State Actualizar segundo momento con gradiente \textbf{actual}: $v \gets \beta_2 v + (1-\beta_2) g_t^2$
    \State $\theta_t \gets \theta_{t-1} - \eta m$
\EndFor
\Return $\theta_T$
\end{algorithmic}
\end{algorithm}

\paragraph{AdEMAMix (Mezcla de Promedios Móviles Exponenciales).}
AdEMAMix reemplaza el EMA único de gradientes de Adam por una \textbf{mezcla de dos EMA}: uno de decadencia rápida y otro de decadencia lenta. Esto aborda la observación de que los gradientes siguen siendo informativos durante decenas de miles de pasos, no solo cientos:
\begin{align}
m_{1,t} &= \beta_1 m_{1,t-1} + (1-\beta_1)g_t \quad\text{(EMA rápido, $\beta_1 \approx 0.9$)} \\
m_{2,t} &= \beta_3 m_{2,t-1} + (1-\beta_3)g_t \quad\text{(EMA lento, $\beta_3 \approx 0.9999$)} \\
v_t &= \beta_2 v_{t-1} + (1-\beta_2)g_t^2 \\
m_t &= m_{1,t} + \alpha_t m_{2,t} \quad\text{(mezcla con peso programado $\alpha_t$)} \\
\theta_{t+1} &= \theta_t - \eta\frac{m_t}{\sqrt{v_t}+\epsilon}
\end{align}
El peso de la mezcla $\alpha_t$ típicamente aumenta durante el entrenamiento, permitiendo que el EMA rápido proporcione adaptación inmediata mientras que el EMA lento acumula correlaciones de gradiente a largo plazo. Empíricamente, un modelo de 1.3B parámetros en 101B tokens alcanza una pérdida similar a AdamW en 197B tokens (ganancia del 95\% en eficiencia de datos), lo que sugiere que el EMA lento reduce significativamente el olvido del modelo.

\begin{algorithm}[H]
\caption{AdEMAMix (Mezcla de Promedios Móviles Exponenciales)}
\label{alg:ademamix}
\begin{algorithmic}[1]
\Require Parámetros iniciales $\theta_0$, tasa de aprendizaje $\eta$, $\beta_1$ (rápido), $\beta_3$ (lento), $\beta_2$ (segundo momento)
\Require Peso de mezcla $\alpha_t$ (típicamente creciente), tolerancia $\epsilon$
\Ensure Parámetros actualizados $\theta_T$
\State Inicializar: $m_1 \gets 0$ (EMA rápido), $m_2 \gets 0$ (EMA lento), $v \gets 0$
\For{$t = 1, 2, \ldots, T$}
    \State $g_t \gets \nabla f(\theta_{t-1})$
    \If{decadencia\_de\_peso $> 0$}
        \State $\theta_{t-1} \gets \theta_{t-1}(1 - \eta\lambda)$
    \EndIf
    \State $m_1 \gets \beta_1 m_1 + (1-\beta_1) g_t$ \Comment{EMA rápido}
    \State $m_2 \gets \beta_3 m_2 + (1-\beta_3) g_t$ \Comment{EMA lento}
    \State $v \gets \beta_2 v + (1-\beta_2) g_t^2$
    \State $m_t \gets m_1 + \alpha_t m_2$ \Comment{Mezcla}
    \State $\theta_t \gets \theta_{t-1} - \eta (m_t / (\sqrt{v} + \epsilon))$
\EndFor
\Return $\theta_T$
\end{algorithmic}
\end{algorithm}

\paragraph{MARS (Haciendo Brillar las Tasas de Aprendizaje Adaptativas).}
MARS combina métodos de gradiente precondicionado (ej., AdamW) con \textbf{reducción de varianza} mediante momentum recursivo estocástico escalado. La innovación central es una estimación del gradiente con varianza reducida:
\begin{align}
c_t &= g_t + \gamma(c_{t-1} - g_{t-1}) \quad\text{(estimación recursiva tipo SVRG)} \\
m_t &= \beta_1 m_{t-1} + (1-\beta_1)c_t \\
v_t &= \beta_2 v_{t-1} + (1-\beta_2)c_t^2 \\
\theta_{t+1} &= \theta_t - \eta\frac{m_t}{\sqrt{v_t}+\epsilon}
\end{align}
donde $\gamma \in [0.01, 0.1]$ controla cuánta información de gradiente histórico se retiene. El gradiente con varianza reducida $c_t$ actúa como un filtro de ruido implícito, amplificando las direcciones de señal consistentes y amortiguando el ruido contradictorio. MARS-AdamW supera consistentemente a AdamW por márgenes significativos en el preentrenamiento de GPT-2, lo que sugiere que la reducción de varianza mejora sustancialmente la convergencia en el entrenamiento estocástico por mini-lotes.

\begin{algorithm}[H]
\caption{MARS (Haciendo Brillar las Tasas de Aprendizaje Adaptativas)}
\label{alg:mars}
\begin{algorithmic}[1]
\Require Parámetros iniciales $\theta_0$, tasa de aprendizaje $\eta$, $\beta_1, \beta_2$, decadencia de peso $\lambda$
\Require Coeficiente de reducción de varianza $\gamma \in [0.01, 0.1]$
\Ensure Parámetros actualizados $\theta_T$
\State Inicializar: $m \gets 0$, $v \gets 0$, $c \gets 0$ (gradiente con varianza reducida), $g_{-1} \gets 0$
\For{$t = 1, 2, \ldots, T$}
    \State $g_t \gets \nabla f(\theta_{t-1})$
    \State Reducción de varianza: $c \gets g_t + \gamma(c - g_{t-1})$
    \State $m \gets \beta_1 m + (1-\beta_1) c$
    \State $v \gets \beta_2 v + (1-\beta_2) c^2$
    \If{decadencia\_de\_peso $> 0$}
        \State $\theta_{t-1} \gets \theta_{t-1}(1 - \eta\lambda)$
    \EndIf
    \State $\theta_t \gets \theta_{t-1} - \eta (m / (\sqrt{v} + \epsilon))$
    \State $g_{t-1} \gets g_t$
\EndFor
\Return $\theta_T$
\end{algorithmic}
\end{algorithm}

\paragraph{Optimizadores Cautelosos (Cautious AdamW, Cautious Lion).}
El marco Cauteloso aplica una \textbf{modificación de una línea} a cualquier optimizador basado en momentum: enmascarar la actualización de modo que solo se apliquen las dimensiones donde el momentum y la dirección del gradiente coinciden:
\begin{align}
\text{mask}_i &= \begin{cases} 1 & \text{si } m_t[i] \cdot g_t[i] > 0 \quad\text{(acuerdo)} \\ 0 & \text{en caso contrario} \end{cases} \\
u_t &= \left(\frac{m_t}{\sqrt{v_t}+\epsilon}\right) \odot \text{mask} \quad\text{(enmascaramiento elemento a elemento)} \\
\theta_{t+1} &= \theta_t - \eta u_t
\end{align}
La intuición: el momentum $m_t$ estima la dirección del gradiente a partir del historial; el gradiente actual $g_t$ es la señal instantánea. Cuando ambos coinciden, el optimizador está seguro y debe actualizar agresivamente. Cuando discrepan, el optimizador está en conflicto---la tendencia histórica apunta hacia una región que pudo haber sido buena antes, pero la evidencia actual la contradice. Al enmascarar las dimensiones en conflicto, Cauteloso se vuelve más conservador y evita pasos corruptos. Empíricamente, este simple enmascaramiento logra hasta \textbf{1.47$\times$ de aceleración} en el preentrenamiento de Llama y MAE mientras preserva las garantías de convergencia.

\begin{algorithm}[H]
\caption{Cautious AdamW (Enmascaramiento de Actualización por Consenso)}
\label{alg:cautious}
\begin{algorithmic}[1]
\Require Parámetros iniciales $\theta_0$, tasa de aprendizaje $\eta$, $\beta_1, \beta_2$, decadencia de peso $\lambda$
\Ensure Parámetros actualizados $\theta_T$
\State Inicializar: $m \gets 0$, $v \gets 0$
\For{$t = 1, 2, \ldots, T$}
    \State $g_t \gets \nabla f(\theta_{t-1})$
    \State $m \gets \beta_1 m + (1-\beta_1) g_t$
    \State $v \gets \beta_2 v + (1-\beta_2) g_t^2$
    \State Máscara de consenso: $\text{mask}_i \gets 1$ si $m[i] \cdot g_t[i] > 0$, si no $0$ \Comment{Acuerdo}
    \State Actualización base Adam: $u \gets m / (\sqrt{v} + \epsilon)$
    \State Aplicar máscara: $u_{\text{masked}} \gets u \odot \text{mask}$ \Comment{Elemento a elemento}
    \If{decadencia\_de\_peso $> 0$}
        \State $\theta_{t-1} \gets \theta_{t-1}(1 - \eta\lambda)$
    \EndIf
    \State $\theta_t \gets \theta_{t-1} - \eta u_{\text{masked}}$
\EndFor
\Return $\theta_T$
\end{algorithmic}
\end{algorithm}

\subsection{Optimizadores de Lote Grande, Eficientes en Memoria y Sin Tasa de Aprendizaje}
\paragraph{LAMB (Layer-wise Adaptive Moments optimizer for Batch training).}
LAMB extiende Adam con \textbf{escalado de tasa adaptativo por capas} (inspirado en LARS), permitiendo entrenamiento estable con tamaños de lote extremos (e.g., 64K en BERT):
\begin{align}
m_t^L &= \beta_1 m_{t-1}^L + (1-\beta_1)g_t^L \quad\text{(momento de primer orden por capas)} \\
v_t^L &= \beta_2 v_{t-1}^L + (1-\beta_2)(g_t^L)^2 \\
u_{\text{adam}}^L &= \frac{m_t^L}{\sqrt{v_t^L}+\epsilon} + \lambda\theta_{t-1}^L \quad\text{(paso adaptativo base con decaimiento)} \\
\phi^L &= \frac{\|\theta_{t-1}^L\|_2}{\|u_{\text{adam}}^L\|_2} \quad\text{(razón de confianza: normalización por capas)} \\
\theta_t^L &= \theta_{t-1}^L - \eta \cdot \phi^L \cdot u_{\text{adam}}^L
\end{align}
La \textbf{razón de confianza} $\phi^L = \|\theta^L\|_2 / \|u_{\text{adam}}^L\|_2$ normaliza el paso de actualización efectivo en relación con la magnitud de los pesos. En lotes muy grandes, el ruido estocástico del gradiente puede superar a la señal, desestabilizando las tasas de aprendizaje por capas. Al escalar las actualizaciones proporcionalmente a las normas de los pesos, LAMB asegura cambios relativos en lugar de absolutos, evitando que actualizaciones pequeñas y seguras en vectores grandes dominen. LAMB preserva los beneficios de generalización de lotes pequeños mientras permite entrenamiento eficiente con lotes grandes.

\begin{algorithm}[H]
\caption{LAMB (Layer-wise Adaptive Moments optimizer for Batch training)}
\label{alg:lamb}
\begin{algorithmic}[1]
\Require Parámetros iniciales $\theta_0$, tasa de aprendizaje $\eta$, $\beta_1, \beta_2$, decaimiento de pesos $\lambda$
\Ensure Parámetros actualizados $\theta_T$
\State Inicializar: Para cada capa $L$: $m^L \gets 0$, $v^L \gets 0$
\For{$t = 1, 2, \ldots, T$}
    \For{cada capa $L$}
        \State $g_t^L \gets \nabla f(\theta_t^L)$
        \State $m^L \gets \beta_1 m^L + (1-\beta_1) g_t^L$
        \State $v^L \gets \beta_2 v^L + (1-\beta_2) (g_t^L)^2$
        \State Adam base: $u_{\text{adam}}^L \gets m^L / (\sqrt{v^L} + \epsilon)$
        \If{weight\_decay $> 0$}
            \State $u_{\text{adam}}^L \gets u_{\text{adam}}^L + \lambda \theta_{t-1}^L$
        \EndIf
        \State Razón de confianza: $\phi^L \gets \|\theta_{t-1}^L\|_2 / \|u_{\text{adam}}^L\|_2$ si ambos son no nulos, sino $1$
        \State $\theta_t^L \gets \theta_{t-1}^L - \eta \phi^L u_{\text{adam}}^L$
    \EndFor
\EndFor
\Return $\theta_T$
\end{algorithmic}
\end{algorithm}

\paragraph{Schedule-Free AdamW.}
Los métodos Schedule-Free eliminan el diseño explícito de programación de la receta de optimización. El algoritmo mantiene dos flujos de parámetros: $z_t$ (exploración) y $x_t$ (promedio suavizado):
\begin{align}
v_t &= \beta_2 v_{t-1} + (1-\beta_2)g_t^2 \quad\text{(varianza)} \\
z_t &= z_{t-1}(1-\eta\lambda) - \eta\frac{g_t}{\sqrt{v_t}+\epsilon} \quad\text{(paso adaptativo con decaimiento)} \\
c_t &= \frac{1}{t+1} \quad\text{(coeficiente de promediado, decrece con el tiempo)} \\
x_t &= (1-c_t)x_{t-1} + c_t z_t \quad\text{(promediado de iteraciones)} \\
y_t &= (1-\beta)z_t + \beta x_t \quad\text{(interpolación para el punto de evaluación)}
\end{align}
El coeficiente de promediado $c_t = 1/(t+1)$ implementa una programación de tasa de aprendizaje implícita sin especificar explícitamente el número total de pasos $T$. Este marco unifica la programación y el promediado de iteraciones: el algoritmo explora mediante $z_t$ mientras acumula dirección estable mediante $x_t$. El punto de evaluación $y_t$ (usado para el cómputo del gradiente) interpola entre exploración y estabilidad. Schedule-Free logra convergencia de última generación en optimización convexa, aprendizaje profundo a gran escala y aprendizaje por refuerzo, eliminando un hiperparámetro importante.

\begin{algorithm}[H]
\caption{Schedule-Free AdamW}
\label{alg:schedulefree}
\begin{algorithmic}[1]
\Require Parámetros iniciales $\theta_0$, tasa de aprendizaje $\eta$ (fija), $\beta_1, \beta_2$, decaimiento de pesos $\lambda$
\Ensure Parámetros actualizados $\theta_T$ o promediado $x_T$
\State Inicializar: $z \gets \theta_0$ (exploración), $x \gets \theta_0$ (promedio), $v \gets 0$, $t \gets 0$
\For{$t = 1, 2, \ldots, T$}
    \State $g_t \gets \nabla f(y_{t-1})$ (gradiente en el punto de interpolación)
    \State $v \gets \beta_2 v + (1-\beta_2) g_t^2$ \Comment{Varianza}
    \State Decaimiento de pesos sobre $z$: $z \gets z(1 - \eta\lambda)$
    \State $z \gets z - \eta g_t / (\sqrt{v} + \epsilon)$ \Comment{Paso adaptativo sobre el gradiente crudo}
    \State $c_t \gets 1 / (t+1)$ \Comment{Coeficiente de promediado}
    \State $x \gets (1-c_t) x + c_t z$ \Comment{Promediado de iteraciones}
    \State $y_t \gets (1-\beta_1) z + \beta_1 x$ \Comment{Interpolación para la siguiente evaluación}
\EndFor
\Return $x_T$ o $y_T$
\end{algorithmic}
\end{algorithm}

\paragraph{Adafactor.}
Adafactor reduce la memoria del optimizador \textbf{factorizando} estadísticas de segundo momento para parámetros con forma de matriz, almacenando solo acumuladores de filas y columnas en lugar de estados de varianza densos. Para una matriz de gradiente $G_t \in \mathbb{R}^{m \times n}$:
\begin{align}
R_t &= \beta_2 R_{t-1} + (1-\beta_2)(G_t^2) \mathbf{1}_n^T \quad\text{(varianza por filas)} \\
C_t &= \beta_2 C_{t-1} + (1-\beta_2) \mathbf{1}_m^\top (G_t^2) \quad\text{(varianza por columnas)} \\
\hat{V}_t &= \frac{R_t C_t}{\mathbf{1}_n^\top R_t} \quad\text{(varianza reconstruida mediante producto exterior)} \\
U_t &= \frac{G_t}{\sqrt{\hat{V}_t}+\epsilon} \quad\text{(paso adaptativo normalizado)} \\
\hat{U}_t &= \frac{U_t}{\max(1, \text{RMS}(U_t))} \quad\text{(recorte de estabilidad)}
\end{align}
En lugar de almacenar $m \times n$ valores de segundo momento, Adafactor almacena solo $m + n$ acumuladores, reduciendo la memoria de $O(n_{\text{params}})$ a $O(\sqrt{n_{\text{params}}})$. Esto es más atractivo cuando la VRAM está dominada por el estado del optimizador en lugar de las activaciones. La contrapartida es una expresividad de optimización reducida y posible inestabilidad en algunas tareas.

\begin{algorithm}[H]
\caption{Adafactor (Estadísticas de Segundo Momento Factorizadas)}
\label{alg:adafactor}
\begin{algorithmic}[1]
\Require Matriz de gradiente $G_t \in \mathbb{R}^{m \times n}$, tasa de aprendizaje $\eta$, $\beta_2 \approx 0.999$
\Ensure Parámetros actualizados
\State Inicializar: Acumuladores de fila $R \gets \epsilon \mathbf{1}^T_m$, Columna $C \gets \epsilon \mathbf{1}_n$
\For{$t = 1, 2, \ldots, T$}
    \State $G_t \gets \nabla f(\theta_t)$ (gradiente matricial)
    \State $R \gets \beta_2 R + (1-\beta_2) (G_t^2) \mathbf{1}_n^T$ \Comment{Productos internos por filas}
    \State $C \gets \beta_2 C + (1-\beta_2) \mathbf{1}_m^\top (G_t^2)$ \Comment{Productos internos por columnas}
    \State Varianza reconstruida: $\hat{V} \gets (R \cdot C) / (\mathbf{1}_n^\top R)$ \Comment{Mediante producto exterior}
    \State Paso adaptativo normalizado: $U_t \gets G_t / (\sqrt{\hat{V}} + \epsilon)$
    \State Recorte RMS por fila: $\hat{U}_t \gets U_t / \max(1, \text{RMS}(U_t))$
    \State $\theta_{t+1} \gets \theta_t - \eta \hat{U}_t$
\EndFor
\end{algorithmic}
\end{algorithm}

\paragraph{GaLore (Gradient Low-Rank Projection).}
GaLore proyecta gradientes 2D en un subespacio de bajo rango antes de la optimización, reduciendo la memoria del estado del optimizador mientras mantiene el escalado adaptativo del paso. Para una matriz de parámetros 2D $W \in \mathbb{R}^{m \times n}$:
\begin{align}
U, S, V &= \text{SVD}(G_t) \quad\text{(calcular descomposición singular)} \\
P &\in \mathbb{R}^{n \times r} \quad\text{o} \quad P^\top \in \mathbb{R}^{r \times m} \quad\text{(seleccionar los $r$ vectores singulares superiores)} \\
G_{\text{low}} &= P^\top G_t \quad\text{(proyectar al espacio de bajo rango)} \\
\Delta_{\text{low}} &= \text{Adam}(G_{\text{low}}) \quad\text{(optimizar en el espacio comprimido)} \\
\Delta &= P \Delta_{\text{low}} \quad\text{(reconstruir en el espacio original)}
\end{align}
Al proyectar en un subespacio de rango $r$ (típicamente $r \ll \min(m,n)$), el estado del optimizador se reduce de $O(mn)$ a $O(r(m+n))$ para parámetros 2D. Este enfoque complementario de ahorro de memoria es especialmente relevante cuando el modelo es demasiado grande para el estado completo del optimizador. GaLore funciona sinérgicamente con otras técnicas y ha mostrado resultados empíricos sólidos en modelos de miles de millones de parámetros.

\begin{algorithm}[H]
\caption{GaLore (Gradient Low-Rank Projection)}
\label{alg:galore}
\begin{algorithmic}[1]
\Require Matriz de parámetros 2D $W \in \mathbb{R}^{m \times n}$, tasa de aprendizaje $\eta$, rango $r \ll \min(m,n)$
\Ensure Parámetros actualizados
\State Inicializar: Estado de Adam en el espacio de bajo rango (si es variante de rango 2)
\For{$t = 1, 2, \ldots, T$}
    \State $G_t \gets \nabla f(W_t)$
    \If{$t \mod K = 1$}  \Comment{SVD periódico}
        \State $U_t, S_t, V_t^\top \gets \text{SVD}(G_t)$ (completa o delgada)
        \State Seleccionar proyección: $P \gets U_t[:, :r]$ o $P \gets V_t[:, :r]$
    \EndIf
    \State Proyectar gradiente: $G_{\text{low}} \gets P^\top G_t$ (o $G_t P^\top$ para proyección derecha)
    \State Ejecutar Adam en bajo rango: $\Delta_{\text{low}} \gets \text{Adam}(G_{\text{low}})$
    \State Reconstruir en espacio original: $\Delta \gets P \Delta_{\text{low}}$ (o $\Delta_{\text{low}} P^\top$)
    \State $W_t \gets W_t - \eta \Delta$
\EndFor
\end{algorithmic}
\end{algorithm}

\paragraph{APOLLO, APOLLO-Mini, y Q-APOLLO.}
La familia APOLLO \citep{zhu2025apollosgdlikememoryadamwlevel} parte de la observación de que el denominador elemento a elemento de AdamW puede \textbf{transformarse en una actualización estructurada de la tasa de aprendizaje}. En lugar de almacenar momentos densos para cada entrada de parámetro, APOLLO proyecta un gradiente matricial $G_t \in \mathbb{R}^{m \times n}$ en un subespacio aleatorio compacto y rastrea momentos estilo Adam allí:
\begin{align}
R_t &= P_t G_t \quad\text{o}\quad R_t = G_t P_t^\top \\
M_t^R &= \beta_1 M_{t-1}^R + (1-\beta_1)R_t \\
V_t^R &= \beta_2 V_{t-1}^R + (1-\beta_2)R_t^2 \\
\widetilde{R}_t &= \frac{M_t^R}{\sqrt{V_t^R}+\epsilon}
\end{align}
El estado proyectado \emph{no} se expande de vuelta a una actualización densa de bajo rango como en los métodos basados en SVD. En su lugar, APOLLO estima un tensor de escalado estructurado $S_t$ en el espacio original. En la variante estándar de APOLLO, ese escalado es \textbf{por canales}: cada fila o canal recibe su propia razón de norma. En APOLLO-Mini, el escalado se reduce a un único escalar \textbf{por tensor}, correspondiente al caso extremo de rango 1 descrito en el artículo. La actualización de parámetros resultante es, por tanto, similar a Adam en adaptación pero mucho más cercana a SGD en costo de estado:
\[
W_t \gets (1-\eta\lambda)W_{t-1} - \eta\,\alpha\,(G_t \odot S_t)
\]
donde $\alpha$ es el factor de escala adicional usado para estabilizar variantes altamente comprimidas.

La contribución práctica del artículo es doble. Primero, APOLLO reemplaza el costoso SVD repetido con \textbf{proyección aleatoria gaussiana} actualizada periódicamente, de modo que la carga computacional es multiplicación de matrices ordinaria en lugar de descomposición espectral. Segundo, el optimizador es inusualmente tolerante a la compresión extrema: incluso \textbf{APOLLO-Mini}, que mantiene solo estado auxiliar de rango 1, sigue siendo competitivo o mejor que AdamW en los experimentos de preentrenamiento reportados, mientras se acerca al costo de memoria de SGD.

\begin{algorithm}[H]
\caption{APOLLO / APOLLO-Mini Escalado de Gradiente Estructurado}
\label{alg:apollo-family}
\begin{algorithmic}[1]
\Require Matriz de pesos $W \in \mathbb{R}^{m \times n}$ con $m \le n$, tasa de aprendizaje $\eta$, factor de escala $\alpha$, tasas de decaimiento $(\beta_1, \beta_2)$, decaimiento de pesos $\lambda$, rango $r$, intervalo de actualización de proyección $T$
\Ensure Parámetros actualizados $W_T$
\State Inicializar momentos proyectados: $M^R \gets 0$, $V^R \gets 0$, paso $t \gets 0$
\Repeat
    \State Calcular gradiente: $G_t \gets \nabla_W \phi(W_t)$
    \If{$t \bmod T = 0$}
        \State Muestrear proyector gaussiano $P_t \sim \mathcal{N}(0, 1/r)$ con una semilla nueva
    \EndIf
    \State Proyectar gradiente: $R_t \gets P_t G_t$ \Comment{o $G_t P_t^\top$ según la disposición}
    \State Actualizar momentos proyectados AdamW: $M_t^R, V_t^R \gets \text{AdamWState}(R_t; \beta_1, \beta_2)$
    \State Normalizar estado proyectado: $\widetilde{R}_t \gets M_t^R / (\sqrt{V_t^R}+\epsilon)$
    \If{APOLLO}
        \State $S_t \gets \operatorname{diag}(s_1^R, \ldots, s_m^R)$, donde $s_i^R = \|\widetilde{R}_t[i,:]\|_2 / \|R_t[i,:]\|_2$
    \Else
        \State $S_t \gets s_t^R$, donde $s_t^R = \|\widetilde{R}_t\|_2 / \|R_t\|_2$
    \EndIf
    \State Actualizar pesos: $W_t \gets (1-\eta\lambda)W_{t-1} - \eta\,\alpha\,(G_t \odot S_t)$
    \State $t \gets t+1$
\Until{convergencia}
\end{algorithmic}
\end{algorithm}

\paragraph{APOLLO-Mini.}
APOLLO-Mini es el miembro de la familia optimizado para \textbf{eficiencia de memoria extrema}. El artículo lo motiva argumentando que, en un espacio compacto de rango 1, el escalado por canales se vuelve demasiado ruidoso, por lo que la actualización se simplifica aún más a un único escalar por tensor. La pérdida de granularidad se compensa parcialmente con el factor de escala explícito $\alpha$ (el artículo discute valores como $128$ en este régimen), proporcionando un punto útil en la frontera de Pareto de optimización: estado muy pequeño, sin sobrecarga de SVD, y aún un comportamiento sólido en preentrenamiento.

\paragraph{Q-APOLLO.}
El artículo de APOLLO también enfatiza que la familia se combina naturalmente con \textbf{cuantización} para entrenamiento con memoria ultra-baja. Este repositorio convierte esa observación de sistemas en una variante concreta del optimizador, \texttt{q\_apollo}. En la implementación, los momentos de primer y segundo orden de bajo rango de APOLLO se almacenan en forma cuantizada junto con escalas y desplazamientos por tensor, y se descuantizan solo cuando es necesario para el siguiente paso. Q-APOLLO preserva, por tanto, la lógica de gradiente proyectado de APOLLO mientras reduce la precisión del estado restante del optimizador, convirtiéndolo en el miembro más agresivo en ahorro de memoria de la pila local de optimizadores.

\paragraph{Anon (Adaptivity Non-restricted Optimizer with Novel convergence technique).}
Anon \citep{anon2026anon} introduce adaptividad ajustable continuamente mediante un único parámetro $\gamma \in \mathbb{R}$, cerrando la brecha entre comportamientos tipo SGD ($\gamma=0$) y tipo Adam ($\gamma=1$), y extrapolando más allá de ambos. La innovación central es el mecanismo Incremental Delay Update (IDU), que permite convergencia estable a través de todo el espectro de adaptividad de valores reales:

\begin{align}
m_t &= \beta_1 m_{t-1} + (1-\beta_1) g_t \quad\text{(momento de primer orden)}\\
v_t &= \beta_2 v_{t-1} + (1-\beta_2) g_t^2 \quad\text{(momento de segundo orden)}\\
\bar{v}_t &= \frac{v_t}{1-\beta_2^{\max(t,1)}} + \epsilon \quad\text{(momento de segundo orden con corrección de sesgo)}\\
\text{fixed\_lr}_k &= \begin{cases}
\bar{v}_k^{-\gamma/2} & \text{si } k = 0 \\
\sqrt{2 / \left(\frac{1}{\text{fixed\_lr}_{k-1}^2} + \bar{v}_k^\gamma\right)} & \text{si } k > 0
\end{cases}\\
\theta_{t+1} &= \theta_t - \frac{\eta}{1-\beta_1^t} \cdot m_t \odot \text{fixed\_lr}_{\lfloor\log_2 t\rfloor}
\end{align}

A diferencia de los optimizadores adaptativos estándar, IDU actualiza el precondicionador solo en potencias de dos ($t = 2^k$), acumulando estadísticas de gradiente entre actualizaciones y agregándolas recursivamente. Este acumulador suave y multiescala evita el seguimiento estricto del máximo de AMSGrad mientras estabiliza demostrablemente la convergencia para cualquier $\gamma \in \mathbb{R}$. El parámetro $\gamma$ controla directamente la adaptividad: $\gamma > 1$ acelera la salida de puntos de silla (favorable para arquitecturas complejas), $\gamma = 1$ recupera el comportamiento tipo Adam, $\gamma = 0$ produce escalado tipo SGD, y $\gamma < 0$ se sesga hacia mínimos más planos (beneficioso para arquitecturas clásicas como CNNs).

Anon mantiene tres buffers de estado ($m$, $v$, fixed\_lr$) con costo $O(n)$ por paso, situándolo en la misma clase de complejidad que AdamW pero con la flexibilidad adicional de adaptividad ajustable. La técnica IDU converge demostrablemente con regret $O(\sqrt{T})$ para problemas convexos y $O(\ln T / \sqrt{T})$ para entornos no convexos, igualando las garantías de los optimizadores adaptativos convencionales mientras soporta todo el rango de adaptividad de valores reales.

\begin{algorithm}[H]
\caption{Anon (Adaptivity Non-restricted Optimizer con IDU)}
\label{alg:anon}
\begin{algorithmic}[1]
\Require Parámetros iniciales $\theta_0$, tasa de aprendizaje $\eta$, $\beta_1, \beta_2$, tolerancia $\epsilon$, adaptividad $\gamma \in \mathbb{R}$
\Ensure Parámetros actualizados $\theta_T$
\State Inicializar: $m \gets 0$, $v \gets 0$, fixed\_lr $\gets 0$, paso $t \gets 0$, contador IDU $k \gets -1$
\For{$t = 1, 2, \ldots, T$}
    \State $g_t \gets \nabla f(\theta_{t-1})$
    \State $m \gets \beta_1 m + (1-\beta_1) g_t$
    \State $v \gets \beta_2 v + (1-\beta_2) g_t^2$
    \If{$t = 2^k$}
        \State $k \gets k+1$
        \State $\bar{v} \gets v/(1 - \beta_2^{\max(t/2,1)}) + \epsilon$
        \If{$k = 0$}
            \State fixed\_lr $\gets \bar{v}^{-\gamma/2}$
        \Else
            \State fixed\_lr $\gets \sqrt{2 / (1/\text{fixed\_lr}^2 + \bar{v}^\gamma)}$
        \EndIf
        \State $v \gets 0$
    \EndIf
    \State $\theta_t \gets \theta_{t-1} - \frac{\eta}{1-\beta_1^t} \cdot m \odot \text{fixed\_lr}$
\EndFor
\Return $\theta_T$
\end{algorithmic}
\end{algorithm}

\paragraph{Prodigy (Approximating the Distance Estimate).}
Prodigy adapta la escala efectiva del paso mediante una estadística tipo distancia acumulativa, eliminando la necesidad de ajuste explícito de la tasa de aprendizaje. El algoritmo mantiene una estimación de distancia acumulativa:
\begin{align}
u_t &= \frac{g_t}{\sqrt{v_t}+\epsilon} \quad\text{(paso adaptativo no normalizado)} \\
s_t &= s_{t-1} + \langle u_t, \theta_t - \theta_0 \rangle \quad\text{(distancia signada acumulativa)} \\
d_t &= \max(d_{t-1}, d_0 + d_{\text{coef}} \cdot s_t) \quad\text{(estimación de distancia con cota inferior)} \\
\theta_{t+1} &= \theta_t - \eta \cdot d_t \cdot u_t
\end{align}
donde $d_0$ es una cota inicial y $d_{\text{coef}}$ controla qué tan agresivamente se adapta la estimación. La estimación de distancia $d_t$ captura la magnitud combinada de gradientes pasados ponderada por el desplazamiento real de los parámetros, implementando una tasa de aprendizaje efectiva implícita. Este escalado consciente de la distancia logra convergencia robusta en distintas escalas de problemas y tamaños de lote sin requerir una programación de la tasa de aprendizaje. Prodigy reduce la sensibilidad a hiperparámetros estimando los tamaños de paso a partir de la geometría de optimización en lugar del conocimiento previo específico del problema.

\begin{algorithm}[H]
\caption{Prodigy (Approximating the Distance Estimate)}
\label{alg:prodigy}
\begin{algorithmic}[1]
\Require Parámetros iniciales $\theta_0$, tasa de aprendizaje $\eta$, coeficiente $d_{\text{coef}}$, distancia inicial $d_0 > 0$
\Ensure Parámetros actualizados $\theta_T$
\State Inicializar: $s \gets 0$ (distancia signada acumulativa), $d \gets d_0$
\For{$t = 1, 2, \ldots, T$}
    \State $g_t \gets \nabla f(\theta_{t-1})$
    \State $v_t \gets \beta_2 v_{t-1} + (1-\beta_2) g_t^2$ \Comment{Segundo momento}
    \State $u_t \gets g_t / (\sqrt{v_t} + \epsilon)$ \Comment{Adaptativo no normalizado}
    \State Actualizar distancia: $s \gets s + \langle u_t, \theta_t - \theta_0 \rangle$ \Comment{Distancia signada}
    \State $d \gets \max(d, d_0 + d_{\text{coef}} \cdot s)$ \Comment{Distancia con cota inferior}
    \State $\theta_t \gets \theta_{t-1} - \eta \cdot d \cdot u_t$ \Comment{Actualización escalada por distancia}
\EndFor
\Return $\theta_T$
\end{algorithmic}
\end{algorithm}

\subsection{Optimizadores de Segundo Orden, Geométricos y de Ortogonalidad}

\begin{algorithm}[H]
\caption{Shampoo (Precondicionamiento Matricial mediante Productos de Kronecker)}
\label{alg:shampoo}
\begin{algorithmic}[1]
\Require Parámetros iniciales $\theta_0$, tasa de aprendizaje $\eta$, intervalo de eigendescomposición $K$
\Ensure Parámetros actualizados $\theta_T$
\State Inicializar: $L \gets \epsilon I_m$, $R \gets \epsilon I_n$ (matrices de Gram izquierda y derecha)
\For{$t = 1, 2, \ldots, T$}
    \State $G \gets \nabla f(\theta_{t-1})$ \Comment{Gradiente}
    \State Actualizar matrices de Gram: $L \gets L + G G^\top$, $R \gets R + G^\top G$
    \If{$t \mod K == 0$}
        \State $Q_L, \Lambda_L \gets \text{eigh}(L)$ \Comment{Eigendescomposición de $L$}
        \State $Q_R, \Lambda_R \gets \text{eigh}(R)$ \Comment{Eigendescomposición de $R$}
        \State $L^{-1/4} \gets Q_L (\Lambda_L + \epsilon)^{-1/4} Q_L^\top$
        \State $R^{-1/4} \gets Q_R (\Lambda_R + \epsilon)^{-1/4} Q_R^\top$
    \EndIf
    \State $\Delta \theta \gets L^{-1/4} G R^{-1/4}$ \Comment{Gradiente precondicionado}
    \State $\theta_t \gets \theta_{t-1} - \eta \cdot \Delta \theta$
\EndFor
\Return $\theta_T$
\end{algorithmic}
\end{algorithm}

\paragraph{Shampoo (Precondicionamiento Matricial mediante Productos de Kronecker).}
Shampoo es un \textbf{método estructurado de segundo orden} que calcula precondicionadores matriciales a partir de estadísticas de producto exterior con estructura de Kronecker. Para una matriz de parámetros 2D $W \in \mathbb{R}^{m \times n}$:
\begin{align}
L_t &= L_{t-1} + G_t G_t^\top \quad\text{(matriz de Gram izquierda/filas, $m \times m$)} \\
R_t &= R_{t-1} + G_t^\top G_t \quad\text{(matriz de Gram derecha/columnas, $n \times n$)} \\
L_t^{-1/4} &= Q_L (\Lambda_L)^{-1/4} Q_L^\top \quad\text{(mediante eigendescomposición)} \\
R_t^{-1/4} &= Q_R (\Lambda_R)^{-1/4} Q_R^\top \\
\Delta W_t &= L_t^{-1/4} G_t R_t^{-1/4} \quad\text{(gradiente precondicionado)}
\end{align}
Shampoo aproxima el precondicionador de matriz completa de Adagrad $H^{-1/2}$ (donde $H$ es el Hessiano) utilizando estructura factorizada de Kronecker. En lugar de almacenar un precondicionador de $(mn) \times (mn)$, Shampoo mantiene dos matrices más pequeñas ($m \times m$ y $n \times n$), capturando correlaciones entre parámetros dentro y entre grupos. El método ha demostrado ser efectivo a escala (sistemas de producción de Google) y es adecuado para arquitecturas transformer con estructura de alto rango. Las eigendescomposiciones se realizan periódicamente (e.g., cada $K=10$ pasos) para amortizar el costo. Contrapartida: mayor cómputo por paso y eigendescomposición $O(m^3 + n^3)$ periódica frente a mejor condicionamiento y convergencia acelerada.

\paragraph{SOAP (Shampoo with Adam in eigenbasis).}
SOAP es una variante simplificada de Shampoo que desacopla el precondicionamiento del seguimiento de momento. En lugar de álgebra matricial compleja sobre gradientes precondicionados, SOAP ejecuta Adam estándar en la base de eigenvectores de los precondicionadores de Shampoo:
\begin{align}
L_t &= L_{t-1} + G_t G_t^\top, \quad R_t = R_{t-1} + G_t^\top G_t \quad\text{(acumulación de Gram)} \\
Q_L, \Lambda_L &= \text{eigh}(L_t), \quad Q_R, \Lambda_R = \text{eigh}(R_t) \quad\text{(eigendescomposición periódica)} \\
G_t^{\text{rot}} &= Q_L^\top G_t Q_R \quad\text{(rotar gradiente a la base de eigenvectores)} \\
m_t^{\text{rot}} &= \beta_1 m_{t-1}^{\text{rot}} + (1-\beta_1) G_t^{\text{rot}} \\
v_t^{\text{rot}} &= \beta_2 v_{t-1}^{\text{rot}} + (1-\beta_2)(G_t^{\text{rot}})^2 \quad\text{(Adam en el espacio rotado)} \\
U_t^{\text{rot}} &= \frac{m_t^{\text{rot}}}{\sqrt{v_t^{\text{rot}}}+\epsilon}, \quad U_t = Q_L U_t^{\text{rot}} Q_R^\top \quad\text{(rotar de vuelta)}
\end{align}
La idea clave: todo el seguimiento de momento ocurre en la base de eigenvectores bien condicionada, simplificando la estabilidad numérica y el análisis teórico. SOAP combina los beneficios de Shampoo (estructura de curvatura explícita) y Adam (mecánica de momento probada), siendo más limpio y a menudo más estable que Shampoo completo. Las eigendescomposiciones se recalculan cada $K$ pasos, amortizando el costo.

\begin{algorithm}[H]
\caption{SOAP (Shampoo with Adam in Eigenbasis)}
\label{alg:soap}
\begin{algorithmic}[1]
\Require Parámetros iniciales $\theta_0$, tasa de aprendizaje $\eta$, $\beta_1, \beta_2$, intervalo de eigendescomposición $K$
\Ensure Parámetros actualizados $\theta_T$
\State Inicializar: $L \gets \epsilon I_m$, $R \gets \epsilon I_n$, $m \gets 0$, $v \gets 0$
\For{$t = 1, 2, \ldots, T$}
    \State $G \gets \nabla f(\theta_{t-1})$ \Comment{Gradiente}
    \State Actualizar Gram: $L \gets L + G G^\top$, $R \gets R + G^\top G$
    \If{$t \mod K == 0$}
        \State $Q_L, \Lambda_L \gets \text{eigh}(L)$, $Q_R, \Lambda_R \gets \text{eigh}(R)$
    \EndIf
    \State $G^{\text{rot}} \gets Q_L^\top G Q_R$ \Comment{Rotar gradiente a la base de eigenvectores}
    \State $m \gets \beta_1 m + (1-\beta_1) G^{\text{rot}}$ \Comment{Media móvil exponencial (corregir sesgo externamente)}
    \State $v \gets \beta_2 v + (1-\beta_2) (G^{\text{rot}})^2$ \Comment{Segundo momento (corregir sesgo externamente)}
    \State $u \gets m / (\sqrt{v} + \epsilon)$ \Comment{Paso adaptativo en la base de eigenvectores}
    \State $\Delta \theta \gets Q_L u Q_R^\top$ \Comment{Rotar de vuelta al espacio de parámetros}
    \State $\theta_t \gets \theta_{t-1} - \eta \cdot \Delta \theta$
\EndFor
\Return $\theta_T$
\end{algorithmic}
\end{algorithm}

\paragraph{Lion (EvoLved Sign Momentum).}
Lion logra un \textbf{sobrecosto de memoria mínimo} al usar actualizaciones basadas en signo en lugar de escalado adaptativo completo:\begin{align}
c_t &= \beta_1 m_t + (1-\beta_1) g_t \quad\text{(entrada de momento)} \\
\theta_{t+1} &= \theta_t - \eta \left(\text{sign}(c_t) + \lambda\theta_t\right) \quad\text{(actualización con operador signo)} \\
m_{t+1} &= \beta_2 m_t + (1-\beta_2) g_t \quad\text{(acumulación de momento)}
\end{align}
Todas las operaciones de escalado elemento a elemento se reemplazan con \texttt{sign()}, que produce $\{-1, 0, +1\}$. Esto reduce drásticamente la memoria en comparación con métodos tipo Adam: Lion almacena solo el buffer de momento, sin varianza de segundo momento. La contrapartida: las actualizaciones basadas en signo sacrifican la adaptación de tasa de aprendizaje elemento a elemento que hace efectivo a Adam. Lion se posiciona no como un optimizador universalmente superior sino como una alternativa especializada de baja memoria y alto rendimiento para escenarios donde la memoria de activaciones domina y se puede sacrificar algo de sofisticación del optimizador. Funciona bien en regímenes de lotes grandes y cuando el rendimiento del hardware es la restricción principal.

\begin{algorithm}[H]
\caption{Lion (EvoLved Sign Momentum)}
\label{alg:lion}
\begin{algorithmic}[1]
\Require Parámetros iniciales $\theta_0$, tasa de aprendizaje $\eta$, $\beta_1, \beta_2$, decaimiento de pesos $\lambda$
\Ensure Parámetros actualizados $\theta_T$
\State Inicializar: $m \gets 0$ (buffer de momento)
\For{$t = 1, 2, \ldots, T$}
    \State $g \gets \nabla f(\theta_{t-1})$ \Comment{Gradiente}
    \State $c \gets \beta_1 m + (1-\beta_1) g$ \Comment{Entrada de momento}
    \State $\theta_t \gets \theta_{t-1} - \eta (\text{sign}(c) + \lambda \theta_{t-1})$ \Comment{Actualización basada en signo con decaimiento de pesos}
    \State $m \gets \beta_2 m + (1-\beta_2) g$ \Comment{Acumulación de momento (desacoplada)}
\EndFor
\Return $\theta_T$
\end{algorithmic}
\end{algorithm}

\paragraph{Sophia (Second-Order Hessian Information with Optimized Approximation).}
Sophia utiliza \textbf{estimaciones diagonales del Hessiano} para escalado consciente de la curvatura sin el costo de memoria y cómputo de los métodos densos de segundo orden.\begin{align}
m_t &= \beta_1 m_{t-1} + (1-\beta_1) g_t \quad\text{(momento de primer orden)} \\
h_t &= \beta_2 h_{t-(k-1)} + (1-\beta_2) \hat{h}_t \quad\text{(Hessiano diagonal, actualizado cada $k$ pasos)} \\
\text{clip}(x, C) &= \min(\max(x, -C), C) \quad\text{(recorte elemento a elemento)} \\
\theta_{t+1} &= \theta_t - \eta \cdot \text{clip}\left(\frac{m_t}{\max(\gamma h_t, \epsilon)}, 1\right)
\end{align}
donde $\hat{h}_t$ es una estimación diagonal (e.g., estimador de traza de Hutchinson) y $\gamma$ es un factor de escala. El Hessiano diagonal $h_t$ captura la curvatura local, permitiendo que el optimizador tome pasos más pequeños en direcciones pronunciadas y pasos más grandes en direcciones planas. La operación de recorte $\text{clip}(\cdot, 1)$ evita que los pasos adaptativos exploten. Sophia pertenece a la familia de métodos conscientes de la curvatura que buscan mejor condicionamiento sin el costo $O(n^2)$ o $O(n^3)$ de los métodos completos de segundo orden. Funciona particularmente bien en escenarios de ajuste fino de segunda pasada.

\begin{algorithm}[H]
\caption{Sophia (Hessiano Diagonal con Actualizaciones Recortadas)}
\label{alg:sophia}
\begin{algorithmic}[1]
\Require Parámetros iniciales $\theta_0$, tasa de aprendizaje $\eta$, $\beta_1, \beta_2$, frecuencia de actualización del Hessiano $k$, umbral de recorte $C$ 
\Ensure Parámetros actualizados $\theta_T$
\State Inicializar: $m \gets 0$, $h \gets \epsilon$ (estimación diagonal del Hessiano)
\For{$t = 1, 2, \ldots, T$}
    \State $g \gets \nabla f(\theta_{t-1})$
    \State $m \gets \beta_1 m + (1-\beta_1) g$ \Comment{Momento de primer orden (corregir sesgo externamente)}
    \If{$t \mod k == 0$}
        \State $\hat{h} \gets \text{HutchinsonEstimate}(g)$ \Comment{Hessiano diagonal mediante Hutchinson}
        \State $h \gets \beta_2 h + (1-\beta_2) \hat{h}$ \Comment{Actualizar estimación del Hessiano}
    \EndIf
    \State $u \gets \frac{m}{\max(\gamma h, \epsilon)}$ \Comment{Paso adaptativo (elemento a elemento)}
    \State $u \gets \text{clip}(u, C)$ \Comment{Recorte elemento a elemento a $[-C, C]$}
    \State $\theta_t \gets \theta_{t-1} - \eta \cdot u$
\EndFor
\Return $\theta_T$
\end{algorithmic}
\end{algorithm}

\paragraph{Muon y Turbo-Muon (Optimizadores Basados en Ortogonalidad).}
Muon y Turbo-Muon son \textbf{optimizadores orientados a la ortogonalidad} que reformulan la geometría de actualización utilizando iteraciones polinomiales de Newton-Schulz para ortogonalización. Para una matriz de parámetros 2D $W$ con gradiente $G_t$:
\begin{align}
X_0 &= \frac{G_t}{\|G_t\|_F + \epsilon} \quad\text{(gradiente normalizado)} \\
A_k &= X_k X_k^\top \quad\text{(Gramiana)} \\
B_k &= b A_k + c A_k^2 \quad\text{(paso polinomial, coeficientes $b, c$ de Newton-Schulz)} \\
X_{k+1} &= a X_k + B_k X_k \quad\text{(iteración $k=0,1,\ldots,4$)} \\
W_{t+1} &= W_t - \eta X_K \quad\text{(actualización con dirección ortogonalizada)}
\end{align}
Muon realiza 5 iteraciones de Newton-Schulz para producir una dirección aproximadamente ortogonal, asegurando que las actualizaciones respeten restricciones geométricas. Turbo-Muon añade un paso de \textbf{precondicionamiento casi-ortogonal} (AOL) antes de las iteraciones, reduciendo el número de pasos de ortogonalización requeridos a 4. La justificación: las actualizaciones ortogonales preservan las normas durante el entrenamiento, evitando la deriva de norma observada en métodos basados en momento. Estos optimizadores muestran potencial en entrenamiento a gran escala pero requieren kernels CUDA personalizados para eficiencia. Contrapartida: alto cómputo por paso (multiplicaciones de matrices e iteraciones polinomiales) frente a una geometría de actualización fundamentalmente mejor condicionada.

\begin{algorithm}[H]
\caption{Muon (Ortogonalización de Newton-Schulz)}
\label{alg:muon}
\begin{algorithmic}[1]
\Require Parámetros iniciales $\theta$ (matrices), tasa de aprendizaje $\eta$, iteraciones Newton-Schulz $K$
\Ensure Parámetros actualizados $\theta$
\For{cada matriz de parámetros 2D $W$}
    \State $G \gets \nabla f(W)$ \Comment{Gradiente}
    \State $X_0 \gets \frac{G}{\|G\|_F + \epsilon}$ \Comment{Normalizar gradiente por norma de Frobenius}
    \For{$k = 0, 1, \ldots, K-1$}
        \State $A_k \gets X_k X_k^\top$ \Comment{Gramiana (costo: $O(mn^2)$ o $O(m^2n)$ para matrices altas/anchas)}
        \State $B_k \gets (3/2) A_k - (1/2) A_k^2$ \Comment{Coeficientes de Newton-Schulz}
        \State $X_{k+1} \gets B_k X_k$ \Comment{Paso polinomial}
    \EndFor
    \State $W \gets W - \eta X_K$ \Comment{Actualizar con dirección ortogonalizada}
\EndFor
\Return $\theta$ actualizado
\end{algorithmic}
\end{algorithm}

\begin{algorithm}[H]
\caption{Turbo-Muon (Ortogonalización Precondicionada con AOL)}
\label{alg:turbo-muon}
\begin{algorithmic}[1]
\Require Parámetros iniciales $\theta$ (matrices), tasa de aprendizaje $\eta$, iteraciones Newton-Schulz $K'$ (típicamente 4)
\Ensure Parámetros actualizados $\theta$
\For{cada matriz de parámetros 2D $W$}
    \State $G \gets \nabla f(W)$ \Comment{Gradiente}
    \State $X_0 \gets \frac{G}{\|G\|_F + \epsilon}$ \Comment{Normalizar gradiente}
    \State \bf{AOL (Almost-Orthogonal preconditioner):} \Comment{Paso de precondicionamiento}
    \State $P \gets (1.5 I - 0.5 X_0 X_0^\top)$ \Comment{Matriz de precondicionamiento}
    \State $X_0 \gets P X_0$ \Comment{Inicio precondicionado}
    \For{$k = 0, 1, \ldots, K'-1$}
        \State $A_k \gets X_k X_k^\top$
        \State $B_k \gets (3/2) A_k - (1/2) A_k^2$
        \State $X_{k+1} \gets B_k X_k$ \Comment{Iteraciones reducidas gracias al precondicionamiento AOL}
    \EndFor
    \State $W \gets W - \eta X_{K'}$ \Comment{Actualizar con dirección ortogonalizada precondicionada}
\EndFor
\Return $\theta$ actualizado
\end{algorithmic}
\end{algorithm}

\small
\begin{longtable}{p{2.5cm}p{3.5cm}p{2.0cm}p{5.3cm}}
\toprule
\textbf{Grupo} & \textbf{Métodos} & \textbf{Objetivo Principal} & \textbf{Interpretación desde el Estudio} \\
\midrule
\endhead
Línea base clásica & SGD, AdamW, RAdam & estabilidad y referencias base & Definen el piso de comparación para las afirmaciones de optimizadores más nuevos. \\
Rediseño de momento & Adan, AdEMAMix, MARS, Cautious AdamW & adaptación de primer orden más rápida o segura & Mejor cuando la velocidad de convergencia o la estabilidad ante gradiente ruidoso es la preocupación principal. \\
Simplificación de lotes grandes y programación & LAMB, Schedule-Free AdamW & robustez operacional a escala & Reducen la fragilidad debida al crecimiento del tamaño de lote o la ingeniería de programación. \\
Eficiencia de memoria & Adafactor, GaLore, APOLLO, APOLLO-Mini, Q-APOLLO, Lion & reducción del estado del optimizador & Más útil cuando la VRAM está dominada por el estado del optimizador en lugar de las activaciones; los métodos de la familia APOLLO reemplazan los momentos densos de AdamW con escalado estructurado proyectado o cuantizado. \\
Consciente de curvatura & Shampoo, SOAP, Sophia & mejor condicionamiento & Preferir cuando una geometría más rica justifica la sobrecarga de implementación y cómputo. \\
Orientado a geometría & Muon, Turbo-Muon & estructura de actualización ortogonalizada & Opciones especializadas para geometría matricial y modelado de representaciones. \\
\bottomrule
\end{longtable}
\normalsize

